\documentclass[11pt]{article}
\usepackage[margin=0.7in]{geometry}
\usepackage{amsmath,amssymb}
\usepackage{enumitem}
\usepackage{tikz}
\usepackage{pgfplots}
\pgfplotsset{compat=1.18}
\usepackage{fancyhdr}
\usepackage{parskip}
\setlist[enumerate]{leftmargin=*, itemsep=6pt}
\pagestyle{fancy}
\fancyhf{}
\fancyfoot[L]{\small\textit{en}Vision\textsuperscript{\textregistered} Algebra 2}
\fancyfoot[C]{\small 7-3 Lesson Quiz, short form}
\fancyfoot[R]{\small\thepage}
\renewcommand{\headrulewidth}{0pt}
\newcommand{\choice}[1]{\tikz[baseline=-0.6ex]\node[draw, circle, inner sep=0.6pt, minimum size=3.3mm, font=\scriptsize\bfseries] {#1};\;}
\newcommand{\checkbox}{$\square$\;}
\newcommand{\tile}[1]{\tikz[baseline=-0.6ex]\node[draw, rounded corners=2pt, minimum width=1.2cm, minimum height=0.7cm] {#1};}

% Short form of the 7-3 Lesson Quiz (a2_7-3_TE.tex, "7-3 Lesson Quiz" block), re-cut 2026-10-31
% with the packet (kept Examples 3 and 4): items 1 (Example 3), 2 (Example 4) and 4 (Example 3B).
% Items 3 (Example 6) and 5 (Example 5) are left off. No item on this form has a figure.

\begin{document}
\relpenalty=10000 \binoppenalty=10000 % keep each inline formula on one line

\noindent\textbf{Name}\; \rule{6cm}{0.4pt} \hfill Period \rule{1.2cm}{0.4pt} \quad Date \rule{2.2cm}{0.4pt}

\vspace{4pt}\noindent\rule{\textwidth}{0.6pt}

\vspace{2pt}\noindent{\Large\textbf{7-3 \; Lesson Quiz: Graphing Sine and Cosine Functions}} \hfill {\small short form}

\vspace{4pt}\noindent\rule{\textwidth}{0.6pt}

{\small Question numbers match the Savvas Lesson Quiz; questions 3 and 5 are not on this form.}

\vspace{4pt}
\textbf{1.} Select \textbf{all} the correct statements.

\vspace{2pt}
\begin{itemize}[leftmargin=1.2em, itemsep=7pt]
\item[] \checkbox \textbf{A.} The amplitude of $y=2\cos\!\left(\frac12x\right)$ is $\frac12$.
\item[] \checkbox \textbf{B.} The period of $y=2\cos\!\left(\frac12x\right)$ is $\pi$.
\item[] \checkbox \textbf{C.} The amplitude of $y=\frac{\pi}{2}\sin\!\left(\frac{\pi}{2}x\right)$ is $\frac{\pi}{2}$.
\item[] \checkbox \textbf{D.} The period of $y=4\pi\sin\!\left(\frac12x\right)$ is $4\pi$.
\item[] \checkbox \textbf{E.} The period of $y=\cos(2\pi x)$ is 1.
\item[] \checkbox \textbf{F.} The amplitude of $y=2\sin\!\left(\frac12x\right)$ is 2.
\end{itemize}

\vspace{4pt}
Show one period calculation: \rule{8cm}{0.4pt}

\vspace{10pt}\noindent\rule{\textwidth}{0.4pt}

\vspace{6pt}
\textbf{2.} What is the average rate of change of the function $y=3\sin(2x)$ on the interval $\left[0,\frac{\pi}{4}\right]$?

\vspace{4pt}
\begin{itemize}[leftmargin=1.2em, itemsep=6pt]
\item[] \choice{A} $-\dfrac{12}{\pi}$ \qquad\qquad \choice{B} $-\dfrac{6}{\pi}$ \qquad\qquad \choice{C} $\dfrac{6}{\pi}$ \qquad\qquad \choice{D} $\dfrac{12}{\pi}$
\end{itemize}

\vspace{4pt}
Show your work: the values of $y$ at $x=0$ and at $x=\frac{\pi}{4}$, and the rate of change.

\vspace{4pt}
\rule{\textwidth}{0.4pt}

\vspace{14pt}
\rule{\textwidth}{0.4pt}

\vspace{10pt}\noindent\rule{\textwidth}{0.4pt}

\vspace{6pt}
\textbf{4.} The graph of the function $f(x)=5\cos(2x)$ is a reflection over the $x$-axis of the graph of
\[ g(x)=\;\rule{1.5cm}{0.4pt}\;\;\rule{1.5cm}{0.4pt}\left(\rule{1.5cm}{0.4pt}\right). \]
Write one tile in each blank.\quad \tile{$-5$}\; \tile{$\cos$}\; \tile{$(2x)$}\; \tile{$\sin$}\; \tile{$(-2x)$}\; \tile{$5$}

\vfill
\noindent\textbf{Scoring:} Question 1 and Question 2 are 10 points each. Question 4 is a bonus of up to 5 points. Total out of 20.

\newpage
\noindent{\LARGE\textbf{7-3 Lesson Quiz, short form: answer key}}

\vspace{4pt}\noindent\rule{\textwidth}{0.6pt}

\raggedright
\begin{enumerate}
\item[\textbf{1.}] \textbf{C, D, E, F.} Amplitude $=|a|$, period $=\frac{2\pi}{|b|}$. A is false: the amplitude of $2\cos\left(\frac12x\right)$ is 2 ($\frac12$ is $b$, not $a$). B is false: its period is $\frac{2\pi}{1/2}=4\pi$ (B multiplies $2\pi$ by $\frac12$, the Try It 3 common error). C: $\left|\frac{\pi}{2}\right|=\frac{\pi}{2}$. D: $\frac{2\pi}{1/2}=4\pi$ (the $4\pi$ in front is the amplitude and does not change the period). E: $\frac{2\pi}{2\pi}=1$. F: $|2|=2$. Savvas quiz item 1, DOK 1, Skills Review and Practice F90. Taught by Example 3 (Try It 3 is the same two formulas). \textbf{10 points:} 2.5 for each of C, D, E, F checked; subtract 2.5 for each of A, B checked; minimum 0.
\item[\textbf{2.}] \textbf{D, $\frac{12}{\pi}$.} The period of $3\sin(2x)$ is $\frac{2\pi}{2}=\pi$, so the key points are $\frac{\pi}{4}$ apart: $(0,0)$ and the maximum $\left(\frac{\pi}{4},3\right)$. Average rate of change $=\dfrac{3\sin\left(\frac{\pi}{2}\right)-3\sin 0}{\frac{\pi}{4}-0}=\dfrac{3}{\pi/4}=\dfrac{12}{\pi}$. A has the wrong sign (the graph rises from 0 to the maximum); C ($\frac{6}{\pi}$, that is $\frac{3}{\pi/2}$) divides by the half period $\frac{\pi}{2}$ instead of $\frac{\pi}{4}$, and B has that error with the wrong sign. Savvas quiz item 2, DOK 1, F31. Taught by Example 4 (the five key points, a zero then a maximum a quarter period later) and the packet's Practice (the same rate of change on Example 4's graphs); the rate formula is the Topic 1 skill. \textbf{10 points:} 3 for $y(0)=0$ and $y\left(\frac{\pi}{4}\right)=3$, 3 for the quotient set up over $\frac{\pi}{4}-0$, 4 for choice D.
\item[\textbf{4.}] \textbf{$-5$; $\cos$; $(2x)$, so $g(x)=-5\cos(2x)$.} Reflecting over the $x$-axis replaces $y$ with $-y$: $-\bigl(5\cos(2x)\bigr)=-5\cos(2x)$, and reflecting $g$ back gives $f$. Also correct: $-5$; $\cos$; $(-2x)$, because cosine is even, so $-5\cos(-2x)=-5\cos(2x)$; give it full credit. The tiles $5$, $\cos$, $(-2x)$ are a trap: $5\cos(-2x)=5\cos(2x)$ is $f$ itself, not its reflection. Savvas quiz item 4, DOK 1, F90. Taught by Example 3B ($y=-f(x)$ is a reflection over the $x$-axis), rehearsed by Try It 4b ($h(x)=-\sin 2\pi x$) and the packet's extension (Practice 23, Part E). \textbf{Bonus 5:} 3 for $-5$, 1 for $\cos$, 1 for $(2x)$ or $(-2x)$.
\end{enumerate}

\textbf{Items left off this form} (available for a retake or Form B): Savvas item 3, the equation $f(x)=2\cos\left(\frac{\pi}{2}x\right)$ against the graph of $g$ (a sine curve through the origin with maxima at $\left(\frac12,2\right)$ and $\left(\frac52,2\right)$): how do the period and the amplitude compare? (answer B: the amplitudes are both 2, and the period of $f$, $\frac{2\pi}{\pi/2}=4$, is twice the period of $g$, $\frac52-\frac12=2$; DOK 2, F90; the cut Example 6); Savvas item 5, the amplitude of the tide model with high tide 5 ft and low tide 1 ft (answer 2 ft, $\frac{5-1}{2}$; DOK 2, F90; the cut Example 5, and Practice 20 is its DOK 3 cousin). Both are left off because no kept Example teaches them. Item 5 is a one-round whiteboard bridge if the teacher puts it back: amplitude $=\frac{\text{max}-\text{min}}{2}$, midline $=\frac{\text{max}+\text{min}}{2}$.

\textbf{IXL replacement:} SmartScore 70 or above on 9S8 and A7V replaces this quiz with 20/20.

\end{document}
